\documentclass[11pt]{article}
% preamble.tex — shared preamble for the dafny-verus-learn-0 written track.
%
% Usage (every <ID>_problems.tex, <ID>_answers.tex and solutions/<ID>_solutions.tex):
%
%   \documentclass[11pt]{article}
%   % preamble.tex — shared preamble for the dafny-verus-learn-0 written track.
%
% Usage (every <ID>_problems.tex, <ID>_answers.tex and solutions/<ID>_solutions.tex):
%
%   \documentclass[11pt]{article}
%   % preamble.tex — shared preamble for the dafny-verus-learn-0 written track.
%
% Usage (every <ID>_problems.tex, <ID>_answers.tex and solutions/<ID>_solutions.tex):
%
%   \documentclass[11pt]{article}
%   \input{preamble}          % the Makefile always runs latexmk from paper/
%   \begin{document}
%   ...
%   \end{document}
%
% Engine: pdflatex. Dafny and Verus sources are mostly ASCII; the `literate`
% tables below map the common Unicode symbols to math glyphs inside listings.
% For other Unicode, extend a table, or build that file with lualatex
% (`latexmk -lualatex <file>`) after replacing the inputenc/fontenc lines with
% \usepackage{fontspec} and a monospace font with good coverage.
%
% Contents: program-logic macros (Hoare triples, wp, VCs, substitution,
% semantics brackets, typing judgments), inference rules (mathpartir),
% theorem and problem environments, the answer environment, the grader
% margin-note macro, and listings languages `dafny` and `verus` (Python uses
% listings' built-in `Python`), with inline forms \dfy{...}, \vrs{...}, \py{...}.
\usepackage[utf8]{inputenc}
\usepackage[T1]{fontenc}

\usepackage{amsmath}
\usepackage{amssymb}
\usepackage{amsthm}
\usepackage{stmaryrd}         % \llbracket \rrbracket for semantics brackets
\usepackage{mathpartir}       % inference rules: \inferrule, the mathpar environment
\usepackage{listings}
\usepackage{xcolor}
\usepackage[margin=1.1in]{geometry}

% ---------------------------------------------------------------------------
% Program logic.
%   \hoare{P}{c}{Q}       {P} c {Q}           (partial correctness)
%   \thoare{P}{c}{Q}      [P] c [Q]           (total correctness)
%   \provable{P}{c}{Q}    |- {P} c {Q}        (derivable in the proof system)
%   \WP{c}{Q}             wp(c, Q)            (\wp is LaTeX's Weierstrass p, so
%   \WLP{c}{Q}            wlp(c, Q)            the macro names are upper case)
%   \VC{c}{Q}             VC(c, Q)
%   \subst{Q}{x}{e}       Q[e/x]              (e replaces x in Q)
%   \sem{e}               [[e]]               (denotation)
%   \assign, \seqc        :=  and the sequencing semicolon
%   \impl                 ==> as a math arrow (LaTeX's own \iff gives <==>)
% ---------------------------------------------------------------------------
\newcommand{\hoare}[3]{\{\,#1\,\}\;#2\;\{\,#3\,\}}
\newcommand{\thoare}[3]{[\,#1\,]\;#2\;[\,#3\,]}
\newcommand{\provable}[3]{\vdash \hoare{#1}{#2}{#3}}
\newcommand{\WP}[2]{\mathrm{wp}(#1,\,#2)}
\newcommand{\WLP}[2]{\mathrm{wlp}(#1,\,#2)}
\newcommand{\VC}[2]{\mathrm{VC}(#1,\,#2)}
\newcommand{\subst}[3]{#1[#3/#2]}
\newcommand{\sem}[1]{\llbracket #1 \rrbracket}
\newcommand{\assign}{\mathrel{:=}}
\newcommand{\seqc}{\mathbin{;}\,}
\newcommand{\impl}{\Longrightarrow}

% ---------------------------------------------------------------------------
% Judgments.  \judgment{\Gamma}{e}{\tau}  ==>  Gamma |- e : tau
% ---------------------------------------------------------------------------
\newcommand{\judgment}[3]{#1 \vdash #2 : #3}
\newcommand{\judg}[3]{\judgment{#1}{#2}{#3}}          % short alias
\newcommand{\emptyjudgment}[2]{\cdot \vdash #1 : #2}  % empty context
\newcommand{\ctx}{\Gamma}
\newcommand{\turnstile}{\vdash}

% ---------------------------------------------------------------------------
% Theorem environments.
% ---------------------------------------------------------------------------
\theoremstyle{plain}
\newtheorem{theorem}{Theorem}[section]
\newtheorem{lemma}[theorem]{Lemma}
\newtheorem{corollary}[theorem]{Corollary}
\newtheorem{proposition}[theorem]{Proposition}
\theoremstyle{definition}
\newtheorem{definition}[theorem]{Definition}
% Written-track problems are numbered B1, B2, ... on both tracks (P<n> is a
% predict item). Label each one \label{prob:B<n>} and its answer \label{ans:B<n>}.
\newtheorem{problem}{Problem}
\renewcommand{\theproblem}{B\arabic{problem}}
\theoremstyle{remark}
\newtheorem*{remark}{Remark}
\newenvironment{answer}
  {\par\medskip\noindent\textbf{Answer.}\begin{quotation}\noindent}
  {\end{quotation}\medskip}

% ---------------------------------------------------------------------------
% Grader margin notes. Rendered form of the `% GRADER[...][...][...]:` comment
% convention (see template_answers.tex). Use inline where a visible note helps:
%   \gradernote{B2/induction-soundness/2: the invariant is not preserved by the last assignment}
% ---------------------------------------------------------------------------
\definecolor{gradercolor}{RGB}{160, 30, 30}
\newcommand{\gradernote}[1]{\textcolor{gradercolor}{\small\sffamily [GRADER: #1]}}

% ---------------------------------------------------------------------------
% Code listings: Dafny, Verus (Rust inside verus!{...}), Python.
% ---------------------------------------------------------------------------
\definecolor{codekeyword}{RGB}{0, 60, 130}
\definecolor{codecomment}{RGB}{90, 110, 90}
\definecolor{codestring}{RGB}{130, 60, 10}
\definecolor{codehole}{RGB}{180, 0, 0}

\lstdefinelanguage{dafny}{
  sensitive=true,
  morekeywords={
    abstract, allocated, as, assert, assume, bool, break, by, calc, case, char,
    class, codatatype, const, constructor, continue, datatype, decreases, else,
    ensures, exists, expect, export, extends, false, forall, fresh, function,
    ghost, if, import, in, include, int, invariant, is, iterator, label, lemma,
    map, match, method, modifies, modify, module, multiset, nat, new, newtype,
    null, old, opaque, opened, predicate, print, provides, reads, real, refines,
    requires, return, returns, reveal, reveals, seq, set, static, string, then,
    this, trait, true, twostate, type, unchanged, var, while, witness},
  emph={assume},
  emphstyle=\color{codehole}\bfseries,
  comment=[l]{//},
  morecomment=[s]{/*}{*/},
  morestring=[b]",
}

\lstdefinelanguage{verus}{
  sensitive=true,
  morekeywords={
    as, break, const, continue, crate, dyn, else, enum, false, fn, for, if,
    impl, in, let, loop, match, mod, move, mut, pub, ref, return, self, Self,
    static, struct, super, trait, true, type, unsafe, use, where, while,
    spec, proof, exec, open, closed, tracked, ghost, requires, ensures,
    recommends, invariant, decreases, assert, forall, exists, choose, by,
    broadcast, uninterp, axiom, returns, int, nat, Seq, Set, Map, Ghost, Tracked},
  emph={assume, admit},
  emphstyle=\color{codehole}\bfseries,
  comment=[l]{//},
  morecomment=[s]{/*}{*/},
  morestring=[b]",
}

\lstset{
  basicstyle=\ttfamily\small,
  keywordstyle=\color{codekeyword}\bfseries,
  commentstyle=\color{codecomment}\itshape,
  stringstyle=\color{codestring},
  showstringspaces=false,
  columns=fullflexible,
  keepspaces=true,
  breaklines=true,
  frame=leftline,
  framerule=0.6pt,
  xleftmargin=1em,
  aboveskip=0.8em,
  belowskip=0.8em,
  literate=
    {→}{{$\to$}}1        {←}{{$\leftarrow$}}1  {⇒}{{$\Rightarrow$}}1
    {⇔}{{$\Leftrightarrow$}}1                  {↦}{{$\mapsto$}}1
    {∀}{{$\forall$}}1    {∃}{{$\exists$}}1     {λ}{{$\lambda$}}1
    {∧}{{$\wedge$}}1     {∨}{{$\vee$}}1        {¬}{{$\neg$}}1
    {≠}{{$\neq$}}1       {≤}{{$\leq$}}1        {≥}{{$\geq$}}1
    {∈}{{$\in$}}1        {∉}{{$\notin$}}1      {⊆}{{$\subseteq$}}1
    {×}{{$\times$}}1     {·}{{$\cdot$}}1       {∅}{{$\varnothing$}}1
    {ℕ}{{$\mathbb{N}$}}1 {ℤ}{{$\mathbb{Z}$}}1
    {⊢}{{$\vdash$}}1     {⟦}{{$\llbracket$}}1  {⟧}{{$\rrbracket$}}1,
}

% Code environments and inline code.
\lstnewenvironment{dafnycode}{\lstset{language=dafny}}{}
\lstnewenvironment{veruscode}{\lstset{language=verus}}{}
\lstnewenvironment{pycode}{\lstset{language=Python}}{}
\newcommand{\dfy}[1]{\lstinline[language=dafny]{#1}}
\newcommand{\vrs}[1]{\lstinline[language=verus]{#1}}
\newcommand{\py}[1]{\lstinline[language=Python]{#1}}
          % the Makefile always runs latexmk from paper/
%   \begin{document}
%   ...
%   \end{document}
%
% Engine: pdflatex. Dafny and Verus sources are mostly ASCII; the `literate`
% tables below map the common Unicode symbols to math glyphs inside listings.
% For other Unicode, extend a table, or build that file with lualatex
% (`latexmk -lualatex <file>`) after replacing the inputenc/fontenc lines with
% \usepackage{fontspec} and a monospace font with good coverage.
%
% Contents: program-logic macros (Hoare triples, wp, VCs, substitution,
% semantics brackets, typing judgments), inference rules (mathpartir),
% theorem and problem environments, the answer environment, the grader
% margin-note macro, and listings languages `dafny` and `verus` (Python uses
% listings' built-in `Python`), with inline forms \dfy{...}, \vrs{...}, \py{...}.
\usepackage[utf8]{inputenc}
\usepackage[T1]{fontenc}

\usepackage{amsmath}
\usepackage{amssymb}
\usepackage{amsthm}
\usepackage{stmaryrd}         % \llbracket \rrbracket for semantics brackets
\usepackage{mathpartir}       % inference rules: \inferrule, the mathpar environment
\usepackage{listings}
\usepackage{xcolor}
\usepackage[margin=1.1in]{geometry}

% ---------------------------------------------------------------------------
% Program logic.
%   \hoare{P}{c}{Q}       {P} c {Q}           (partial correctness)
%   \thoare{P}{c}{Q}      [P] c [Q]           (total correctness)
%   \provable{P}{c}{Q}    |- {P} c {Q}        (derivable in the proof system)
%   \WP{c}{Q}             wp(c, Q)            (\wp is LaTeX's Weierstrass p, so
%   \WLP{c}{Q}            wlp(c, Q)            the macro names are upper case)
%   \VC{c}{Q}             VC(c, Q)
%   \subst{Q}{x}{e}       Q[e/x]              (e replaces x in Q)
%   \sem{e}               [[e]]               (denotation)
%   \assign, \seqc        :=  and the sequencing semicolon
%   \impl                 ==> as a math arrow (LaTeX's own \iff gives <==>)
% ---------------------------------------------------------------------------
\newcommand{\hoare}[3]{\{\,#1\,\}\;#2\;\{\,#3\,\}}
\newcommand{\thoare}[3]{[\,#1\,]\;#2\;[\,#3\,]}
\newcommand{\provable}[3]{\vdash \hoare{#1}{#2}{#3}}
\newcommand{\WP}[2]{\mathrm{wp}(#1,\,#2)}
\newcommand{\WLP}[2]{\mathrm{wlp}(#1,\,#2)}
\newcommand{\VC}[2]{\mathrm{VC}(#1,\,#2)}
\newcommand{\subst}[3]{#1[#3/#2]}
\newcommand{\sem}[1]{\llbracket #1 \rrbracket}
\newcommand{\assign}{\mathrel{:=}}
\newcommand{\seqc}{\mathbin{;}\,}
\newcommand{\impl}{\Longrightarrow}

% ---------------------------------------------------------------------------
% Judgments.  \judgment{\Gamma}{e}{\tau}  ==>  Gamma |- e : tau
% ---------------------------------------------------------------------------
\newcommand{\judgment}[3]{#1 \vdash #2 : #3}
\newcommand{\judg}[3]{\judgment{#1}{#2}{#3}}          % short alias
\newcommand{\emptyjudgment}[2]{\cdot \vdash #1 : #2}  % empty context
\newcommand{\ctx}{\Gamma}
\newcommand{\turnstile}{\vdash}

% ---------------------------------------------------------------------------
% Theorem environments.
% ---------------------------------------------------------------------------
\theoremstyle{plain}
\newtheorem{theorem}{Theorem}[section]
\newtheorem{lemma}[theorem]{Lemma}
\newtheorem{corollary}[theorem]{Corollary}
\newtheorem{proposition}[theorem]{Proposition}
\theoremstyle{definition}
\newtheorem{definition}[theorem]{Definition}
% Written-track problems are numbered B1, B2, ... on both tracks (P<n> is a
% predict item). Label each one \label{prob:B<n>} and its answer \label{ans:B<n>}.
\newtheorem{problem}{Problem}
\renewcommand{\theproblem}{B\arabic{problem}}
\theoremstyle{remark}
\newtheorem*{remark}{Remark}
\newenvironment{answer}
  {\par\medskip\noindent\textbf{Answer.}\begin{quotation}\noindent}
  {\end{quotation}\medskip}

% ---------------------------------------------------------------------------
% Grader margin notes. Rendered form of the `% GRADER[...][...][...]:` comment
% convention (see template_answers.tex). Use inline where a visible note helps:
%   \gradernote{B2/induction-soundness/2: the invariant is not preserved by the last assignment}
% ---------------------------------------------------------------------------
\definecolor{gradercolor}{RGB}{160, 30, 30}
\newcommand{\gradernote}[1]{\textcolor{gradercolor}{\small\sffamily [GRADER: #1]}}

% ---------------------------------------------------------------------------
% Code listings: Dafny, Verus (Rust inside verus!{...}), Python.
% ---------------------------------------------------------------------------
\definecolor{codekeyword}{RGB}{0, 60, 130}
\definecolor{codecomment}{RGB}{90, 110, 90}
\definecolor{codestring}{RGB}{130, 60, 10}
\definecolor{codehole}{RGB}{180, 0, 0}

\lstdefinelanguage{dafny}{
  sensitive=true,
  morekeywords={
    abstract, allocated, as, assert, assume, bool, break, by, calc, case, char,
    class, codatatype, const, constructor, continue, datatype, decreases, else,
    ensures, exists, expect, export, extends, false, forall, fresh, function,
    ghost, if, import, in, include, int, invariant, is, iterator, label, lemma,
    map, match, method, modifies, modify, module, multiset, nat, new, newtype,
    null, old, opaque, opened, predicate, print, provides, reads, real, refines,
    requires, return, returns, reveal, reveals, seq, set, static, string, then,
    this, trait, true, twostate, type, unchanged, var, while, witness},
  emph={assume},
  emphstyle=\color{codehole}\bfseries,
  comment=[l]{//},
  morecomment=[s]{/*}{*/},
  morestring=[b]",
}

\lstdefinelanguage{verus}{
  sensitive=true,
  morekeywords={
    as, break, const, continue, crate, dyn, else, enum, false, fn, for, if,
    impl, in, let, loop, match, mod, move, mut, pub, ref, return, self, Self,
    static, struct, super, trait, true, type, unsafe, use, where, while,
    spec, proof, exec, open, closed, tracked, ghost, requires, ensures,
    recommends, invariant, decreases, assert, forall, exists, choose, by,
    broadcast, uninterp, axiom, returns, int, nat, Seq, Set, Map, Ghost, Tracked},
  emph={assume, admit},
  emphstyle=\color{codehole}\bfseries,
  comment=[l]{//},
  morecomment=[s]{/*}{*/},
  morestring=[b]",
}

\lstset{
  basicstyle=\ttfamily\small,
  keywordstyle=\color{codekeyword}\bfseries,
  commentstyle=\color{codecomment}\itshape,
  stringstyle=\color{codestring},
  showstringspaces=false,
  columns=fullflexible,
  keepspaces=true,
  breaklines=true,
  frame=leftline,
  framerule=0.6pt,
  xleftmargin=1em,
  aboveskip=0.8em,
  belowskip=0.8em,
  literate=
    {→}{{$\to$}}1        {←}{{$\leftarrow$}}1  {⇒}{{$\Rightarrow$}}1
    {⇔}{{$\Leftrightarrow$}}1                  {↦}{{$\mapsto$}}1
    {∀}{{$\forall$}}1    {∃}{{$\exists$}}1     {λ}{{$\lambda$}}1
    {∧}{{$\wedge$}}1     {∨}{{$\vee$}}1        {¬}{{$\neg$}}1
    {≠}{{$\neq$}}1       {≤}{{$\leq$}}1        {≥}{{$\geq$}}1
    {∈}{{$\in$}}1        {∉}{{$\notin$}}1      {⊆}{{$\subseteq$}}1
    {×}{{$\times$}}1     {·}{{$\cdot$}}1       {∅}{{$\varnothing$}}1
    {ℕ}{{$\mathbb{N}$}}1 {ℤ}{{$\mathbb{Z}$}}1
    {⊢}{{$\vdash$}}1     {⟦}{{$\llbracket$}}1  {⟧}{{$\rrbracket$}}1,
}

% Code environments and inline code.
\lstnewenvironment{dafnycode}{\lstset{language=dafny}}{}
\lstnewenvironment{veruscode}{\lstset{language=verus}}{}
\lstnewenvironment{pycode}{\lstset{language=Python}}{}
\newcommand{\dfy}[1]{\lstinline[language=dafny]{#1}}
\newcommand{\vrs}[1]{\lstinline[language=verus]{#1}}
\newcommand{\py}[1]{\lstinline[language=Python]{#1}}
          % the Makefile always runs latexmk from paper/
%   \begin{document}
%   ...
%   \end{document}
%
% Engine: pdflatex. Dafny and Verus sources are mostly ASCII; the `literate`
% tables below map the common Unicode symbols to math glyphs inside listings.
% For other Unicode, extend a table, or build that file with lualatex
% (`latexmk -lualatex <file>`) after replacing the inputenc/fontenc lines with
% \usepackage{fontspec} and a monospace font with good coverage.
%
% Contents: program-logic macros (Hoare triples, wp, VCs, substitution,
% semantics brackets, typing judgments), inference rules (mathpartir),
% theorem and problem environments, the answer environment, the grader
% margin-note macro, and listings languages `dafny` and `verus` (Python uses
% listings' built-in `Python`), with inline forms \dfy{...}, \vrs{...}, \py{...}.
\usepackage[utf8]{inputenc}
\usepackage[T1]{fontenc}

\usepackage{amsmath}
\usepackage{amssymb}
\usepackage{amsthm}
\usepackage{stmaryrd}         % \llbracket \rrbracket for semantics brackets
\usepackage{mathpartir}       % inference rules: \inferrule, the mathpar environment
\usepackage{listings}
\usepackage{xcolor}
\usepackage[margin=1.1in]{geometry}

% ---------------------------------------------------------------------------
% Program logic.
%   \hoare{P}{c}{Q}       {P} c {Q}           (partial correctness)
%   \thoare{P}{c}{Q}      [P] c [Q]           (total correctness)
%   \provable{P}{c}{Q}    |- {P} c {Q}        (derivable in the proof system)
%   \WP{c}{Q}             wp(c, Q)            (\wp is LaTeX's Weierstrass p, so
%   \WLP{c}{Q}            wlp(c, Q)            the macro names are upper case)
%   \VC{c}{Q}             VC(c, Q)
%   \subst{Q}{x}{e}       Q[e/x]              (e replaces x in Q)
%   \sem{e}               [[e]]               (denotation)
%   \assign, \seqc        :=  and the sequencing semicolon
%   \impl                 ==> as a math arrow (LaTeX's own \iff gives <==>)
% ---------------------------------------------------------------------------
\newcommand{\hoare}[3]{\{\,#1\,\}\;#2\;\{\,#3\,\}}
\newcommand{\thoare}[3]{[\,#1\,]\;#2\;[\,#3\,]}
\newcommand{\provable}[3]{\vdash \hoare{#1}{#2}{#3}}
\newcommand{\WP}[2]{\mathrm{wp}(#1,\,#2)}
\newcommand{\WLP}[2]{\mathrm{wlp}(#1,\,#2)}
\newcommand{\VC}[2]{\mathrm{VC}(#1,\,#2)}
\newcommand{\subst}[3]{#1[#3/#2]}
\newcommand{\sem}[1]{\llbracket #1 \rrbracket}
\newcommand{\assign}{\mathrel{:=}}
\newcommand{\seqc}{\mathbin{;}\,}
\newcommand{\impl}{\Longrightarrow}

% ---------------------------------------------------------------------------
% Judgments.  \judgment{\Gamma}{e}{\tau}  ==>  Gamma |- e : tau
% ---------------------------------------------------------------------------
\newcommand{\judgment}[3]{#1 \vdash #2 : #3}
\newcommand{\judg}[3]{\judgment{#1}{#2}{#3}}          % short alias
\newcommand{\emptyjudgment}[2]{\cdot \vdash #1 : #2}  % empty context
\newcommand{\ctx}{\Gamma}
\newcommand{\turnstile}{\vdash}

% ---------------------------------------------------------------------------
% Theorem environments.
% ---------------------------------------------------------------------------
\theoremstyle{plain}
\newtheorem{theorem}{Theorem}[section]
\newtheorem{lemma}[theorem]{Lemma}
\newtheorem{corollary}[theorem]{Corollary}
\newtheorem{proposition}[theorem]{Proposition}
\theoremstyle{definition}
\newtheorem{definition}[theorem]{Definition}
% Written-track problems are numbered B1, B2, ... on both tracks (P<n> is a
% predict item). Label each one \label{prob:B<n>} and its answer \label{ans:B<n>}.
\newtheorem{problem}{Problem}
\renewcommand{\theproblem}{B\arabic{problem}}
\theoremstyle{remark}
\newtheorem*{remark}{Remark}
\newenvironment{answer}
  {\par\medskip\noindent\textbf{Answer.}\begin{quotation}\noindent}
  {\end{quotation}\medskip}

% ---------------------------------------------------------------------------
% Grader margin notes. Rendered form of the `% GRADER[...][...][...]:` comment
% convention (see template_answers.tex). Use inline where a visible note helps:
%   \gradernote{B2/induction-soundness/2: the invariant is not preserved by the last assignment}
% ---------------------------------------------------------------------------
\definecolor{gradercolor}{RGB}{160, 30, 30}
\newcommand{\gradernote}[1]{\textcolor{gradercolor}{\small\sffamily [GRADER: #1]}}

% ---------------------------------------------------------------------------
% Code listings: Dafny, Verus (Rust inside verus!{...}), Python.
% ---------------------------------------------------------------------------
\definecolor{codekeyword}{RGB}{0, 60, 130}
\definecolor{codecomment}{RGB}{90, 110, 90}
\definecolor{codestring}{RGB}{130, 60, 10}
\definecolor{codehole}{RGB}{180, 0, 0}

\lstdefinelanguage{dafny}{
  sensitive=true,
  morekeywords={
    abstract, allocated, as, assert, assume, bool, break, by, calc, case, char,
    class, codatatype, const, constructor, continue, datatype, decreases, else,
    ensures, exists, expect, export, extends, false, forall, fresh, function,
    ghost, if, import, in, include, int, invariant, is, iterator, label, lemma,
    map, match, method, modifies, modify, module, multiset, nat, new, newtype,
    null, old, opaque, opened, predicate, print, provides, reads, real, refines,
    requires, return, returns, reveal, reveals, seq, set, static, string, then,
    this, trait, true, twostate, type, unchanged, var, while, witness},
  emph={assume},
  emphstyle=\color{codehole}\bfseries,
  comment=[l]{//},
  morecomment=[s]{/*}{*/},
  morestring=[b]",
}

\lstdefinelanguage{verus}{
  sensitive=true,
  morekeywords={
    as, break, const, continue, crate, dyn, else, enum, false, fn, for, if,
    impl, in, let, loop, match, mod, move, mut, pub, ref, return, self, Self,
    static, struct, super, trait, true, type, unsafe, use, where, while,
    spec, proof, exec, open, closed, tracked, ghost, requires, ensures,
    recommends, invariant, decreases, assert, forall, exists, choose, by,
    broadcast, uninterp, axiom, returns, int, nat, Seq, Set, Map, Ghost, Tracked},
  emph={assume, admit},
  emphstyle=\color{codehole}\bfseries,
  comment=[l]{//},
  morecomment=[s]{/*}{*/},
  morestring=[b]",
}

\lstset{
  basicstyle=\ttfamily\small,
  keywordstyle=\color{codekeyword}\bfseries,
  commentstyle=\color{codecomment}\itshape,
  stringstyle=\color{codestring},
  showstringspaces=false,
  columns=fullflexible,
  keepspaces=true,
  breaklines=true,
  frame=leftline,
  framerule=0.6pt,
  xleftmargin=1em,
  aboveskip=0.8em,
  belowskip=0.8em,
  literate=
    {→}{{$\to$}}1        {←}{{$\leftarrow$}}1  {⇒}{{$\Rightarrow$}}1
    {⇔}{{$\Leftrightarrow$}}1                  {↦}{{$\mapsto$}}1
    {∀}{{$\forall$}}1    {∃}{{$\exists$}}1     {λ}{{$\lambda$}}1
    {∧}{{$\wedge$}}1     {∨}{{$\vee$}}1        {¬}{{$\neg$}}1
    {≠}{{$\neq$}}1       {≤}{{$\leq$}}1        {≥}{{$\geq$}}1
    {∈}{{$\in$}}1        {∉}{{$\notin$}}1      {⊆}{{$\subseteq$}}1
    {×}{{$\times$}}1     {·}{{$\cdot$}}1       {∅}{{$\varnothing$}}1
    {ℕ}{{$\mathbb{N}$}}1 {ℤ}{{$\mathbb{Z}$}}1
    {⊢}{{$\vdash$}}1     {⟦}{{$\llbracket$}}1  {⟧}{{$\rrbracket$}}1,
}

% Code environments and inline code.
\lstnewenvironment{dafnycode}{\lstset{language=dafny}}{}
\lstnewenvironment{veruscode}{\lstset{language=verus}}{}
\lstnewenvironment{pycode}{\lstset{language=Python}}{}
\newcommand{\dfy}[1]{\lstinline[language=dafny]{#1}}
\newcommand{\vrs}[1]{\lstinline[language=verus]{#1}}
\newcommand{\py}[1]{\lstinline[language=Python]{#1}}

\title{R16 --- Capstone + Lab $\gamma$ \\ \large Written problems}
\author{dafny-verus-learn-0}
\date{}
\begin{document}
\maketitle
\sloppy
\section*{Instructions}
Commit the R16 predictions before reading B1 (\texttt{scripts/firsttry.sh R16
--check}). B1 is the blueprint: write it after the chapter's ``On paper
first'' section and before you write any statement in
\texttt{R16\_Statements.dfy} or \texttt{statements.rs}; run
\texttt{/grade-paper R16} in a fresh session and get at least 3 on every B1
dimension before you formalize anything (design \S7.1's hard-tier rule). B2
and B3 are paper cards: use exactly the pinned local range, and name missing
information instead of inventing it. B4 reviews candidate statements for the
emitter: answer it once your own H1 statements are written and before your
statement-only review (\texttt{/review R16 --statements}). B5 is Lab
$\gamma$'s pre-registration: it is graded by \texttt{/grade-paper R16}, every
method dimension must be at least 3 before any lab run, and \texttt{fvh}
refuses a run whose PREREG hash is not on an unlocked gradebook row; it also
refuses to unpack the crate tasks before your statement files are committed.
Claude is both the grader of this paper and a subject of Lab $\gamma$; say so
wherever it matters. The one-page memo, contract R-v2 and the final
\texttt{TRUST.md} are the chapter's last section's work, reviewed in the code
and oral review, not here.

% item: R16.B1
\begin{problem}[informal-proof-first]\label{prob:B1}
\textbf{The blueprint} (45 minutes). Layer A5 gives the VM (\dfy{Accepts},
opaque, used through its four one-step lemmas and \dfy{AcceptsInside};
\dfy{AcceptsFrom}, acceptance with some fuel), the emitter (\dfy{Emit(r, pc)},
the code of \dfy{r} placed at address \dfy{pc}, which falls through to the
address after it; \dfy{EmitProg(r)}, that code at address 0 followed by one
\dfy{IMatch}), \dfy{CodeAt} and the layout lemmas of each constructor's code.
The Verus names are \vrs{accepts}, \vrs{accepts_from}, \vrs{emit},
\vrs{clen}, \vrs{prog} and \vrs{code_at}. The capstone's last theorem is H5's,
whose statement your exercise file prints: the compiled program accepts
exactly the words \dfy{r} matches. Before you write any statement, write, in
both languages wherever they differ:
\begin{enumerate}
\item[(a)] whether that whole-program fact can be proved, in each direction,
  by structural induction on \dfy{r} as it stands; if not, the first
  constructor case where the induction hypothesis is not the fact the case
  needs, and what the case needs instead;
\item[(b)] the two statements you will freeze, completeness and soundness, as
  formulas over A5's names, for the code of \dfy{r} at \dfy{pc} inside an
  arbitrary program, with every hypothesis and every quantifier's range; why
  neither statement may say where the program ends; what the soundness
  conclusion must bound; and what more, if anything, the Verus form needs;
\item[(c)] how (b) gives back the whole-program fact in each direction: which
  continuation, which word it accepts, and with what fuel;
\item[(d)] completeness's proof: the fuelled helper you will state (its exact
  \texttt{requires} and \texttt{ensures}), its induction and measure, and each
  constructor's case: the layout lemma and the one-step lemmas it uses, and how
  the fuels of the pieces compose; for \dfy{Star}, what the inner induction
  runs on and why its first chunk must be non-empty;
\item[(e)] soundness's proof: the fuelled helper's statement, its measure,
  and why the measure decreases in every case (name the cases that spend no
  fuel before they recurse, and what they recurse on); how an \dfy{ISplit}'s two
  branches are handled; and how a pass through \dfy{Star}'s body that reads
  nothing is treated;
\item[(f)] jump threading (the file's \dfy{Hop} and \dfy{Shortcut}): the two
  fuelled directions, each as a statement; for each, how the fuel of one run
  bounds the fuel of the other; and the one instruction case where the two
  programs differ;
\item[(g)] every step you cannot yet justify, declared as a gap.
\end{enumerate}
The test of a blueprint: could a contractor formalize it without inventing a
statement? Write formulas in \LaTeX{} or in \dfy{...}/\vrs{...} code; the
grader reads either. Stating the shapes is the task: copy no statement from
anywhere, and write no proof body for H2--H5 here.
\end{problem}

% item: R16.B2
\begin{problem}[paper-card]\label{prob:B2}
Read local card 07.3, \emph{Inductive Deductive Synthesis: Enabling AI to
Generate Formally Verified Systems}. Use exactly lines 2877--2907 of your
local \texttt{cheatsheets.md}; the excerpt ends before 07.4. Verify the local
bytes:
\begin{verbatim}
sed -n '2877,2907p' ~/code/research/agentic-formal-veri/cheatsheets.md | sha256sum
4ce943081668ca8b55f34e2a5fa18422c8c6f0d444bc542c39fb7ed97fe5c75a  -
\end{verbatim}
A mismatch needs an author re-pin. Fill Task shape, Oracle, Trust, Cheat
surface, Metric, Budget, Variance and Artifact; the card score is the minimum
of the eight field scores. Keep every reported figure with its row and column
(system, benchmark, specification set) and do not pool what the card keeps
apart.

Then tie the card to R16, whose record in PLAN reads it as implementation and
proof grown together, with stalled designs discarded. (a) Map the card's step
kinds and its rule for deferred holes onto R16's workflow (the holes the gate
counts, the blueprint, the frozen statements, the proofs), and say which of
the card's mechanisms has no counterpart in a course that freezes its
statements before the proofs and gives the emitter. (b) Name the event in R16
that plays the part of a discarded design: what stalls, how you would know
that it stalled, and what replaces it; compare the card's trigger for
discarding a design with yours. (c) The card names an audit of finished
proofs: say what it checks, quoting the excerpt, and what the excerpt says
about its limits; then give one of B4's candidates whose proof such an audit
could pass although the statement is defective, with the course instrument
that catches it. (d) Say what the card's ablations do and do not license about
R16's own process (one learner, one capstone, no repeats), keeping each figure
with its row.
\end{problem}

% item: R16.B3
\begin{problem}[paper-card]\label{prob:B3}
Read local card 12.9, \emph{Autonomous Research Agents: A Survey of AI
Scientists and the Verification Gap}. Use exactly lines 5418--5445 of your
local \texttt{cheatsheets.md}; the excerpt ends before 12.10. Verify the local
bytes:
\begin{verbatim}
sed -n '5418,5445p' ~/code/research/agentic-formal-veri/cheatsheets.md | sha256sum
ad66a46bbab0a8836a15cefe8ab48b162310515273b479a6b46747d0a1bf9e43  -
\end{verbatim}
Fill the eight fields as in B2. Keep every rate with the number of entries it
is computed over, and attribute to the card, not to the paper, any figure or
notation that the excerpt marks as the card's own.

Then tie the card to R16, whose record in PLAN reads it as the verification
gap: R16 ships a system and a claim that an outside checker can re-run.
(a) R16 makes two claims: \texttt{rexc}'s whole-program theorem and Lab
$\gamma$'s inflation estimate. Place each claim, and each check that supports
it (the gate's verdict, \texttt{rexc trust-report}, the CLI's tests, the
batch's logs), and each of Lab $\gamma$'s naive graders, on the card's
verification ladder, naming the rungs as the excerpt does. (b) Using the
card's autonomy levels, say where Lab $\gamma$ sits as a system and where one
delivery run's loop sits, and why. (c) List what an outside checker needs to
re-run each of the two claims (commands, pins, artifacts) and what that
checker must still trust; say whether the card's ladder accounts for that
trust, quoting the excerpt. (d) Name one item of the card's disclosure
checklist that Lab $\gamma$'s report cannot satisfy as designed, and why.
\end{problem}

% item: R16.B4
\begin{problem}[statement-review]\label{prob:B4}
Six frozen candidate statements were proposed for the emitter's correctness
facts, over layers A1 and A5 (\dfy{Re}, \dfy{Matches}, \dfy{Prog},
\dfy{Inst}, \dfy{CodeAt}, \dfy{Emit}, \dfy{AcceptsFrom}). They are printed
below in Dafny with the bodies omitted; part (d) asks about Verus. The labels
C1--C6 carry no meaning, and a candidate may be true and still be the wrong
statement to freeze.
Do B4 only after your own statements (H1) are committed and frozen: the
candidates are near misses of the statements you write, so reading them first
would answer H1 for you.

\begin{dafnycode}
lemma CandC1R16(q: Prog, e: Re, a: int, x: seq<char>, y: seq<char>)
  requires CodeAt(q, a, Emit(e, a))
  requires |q| < a + |Emit(e, a)|
  requires Matches(e, x)
  ensures AcceptsFrom(q, a, x + y)

lemma CandC2R16(q: Prog, e: Re, a: int, z: seq<char>)
  requires CodeAt(q, a, Emit(e, a))
  requires AcceptsFrom(q, a, z)
  ensures exists x, y :: z == x + y && Matches(e, x)

lemma CandC3R16(q: Prog, e: Re, a: int, x: seq<char>)
  requires CodeAt(q, a, Emit(e, a))
  requires a + |Emit(e, a)| < |q|
  requires q[a + |Emit(e, a)|] == IMatch
  requires Matches(e, x)
  ensures AcceptsFrom(q, a, x)

lemma CandC4R16(q: Prog, e: Re, a: int, x: seq<char>, y: seq<char>)
  requires CodeAt(q, a, Emit(e, a))
  requires Matches(e, x)
  requires AcceptsFrom(q, a + |Emit(e, a)| + 1, y)
  ensures AcceptsFrom(q, a, x + y)

lemma CandC5R16(q: Prog, e: Re, a: int, z: seq<char>)
  requires CodeAt(q, a, Emit(e, a))
  requires AcceptsFrom(q, a, z)
  ensures exists j :: Matches(e, z[..j]) && AcceptsFrom(q, a + |Emit(e, a)|, z[j..])

lemma CandC6R16(q: Prog, e: Re, a: int, x: seq<char>, y: seq<char>)
  requires CodeAt(q, a, Emit(e, a))
  requires a + |Emit(e, a)| == |q|
  requires Matches(e, x)
  requires AcceptsFrom(q, a + |Emit(e, a)|, y)
  ensures AcceptsFrom(q, a, x + y)
\end{dafnycode}

For each candidate: (a) classify it as faithful or name its defect (for
instance vacuous, weak, false, ill-formed, or true but not inductive: no
proof of it by induction on itself goes through; a class may occur more than
once or not at all); (b) give a witness: a concrete program \dfy{q}, regex
\dfy{e}, address \dfy{a} and words, and what the candidate forces, allows,
forbids or fails to say there, against what the emitter's code actually does
(for a faithful candidate, argue why no witness exists); (c) say which
instrument exposes the defect, and what it shows: a faithful proof fails; a
body that simulates nothing verifies; the verifier rejects the statement
itself; a caller cannot call it, or cannot conclude what it needs from it; or
an induction cannot use it as its own hypothesis. (d) For C5, say what the
same quantifier means in Verus, where the slices are written with
\vrs{subrange}, what the verifier then reports, and which instrument exposes
the defect there. (e) State what a verified body says about a statement.
Write the witnesses as values you could run; write no body, invariant or proof
for H2--H5 here.
\end{problem}

% item: R16.B5
\begin{problem}[pre-registration]\label{prob:B5}
Pre-register Lab $\gamma$ (PLAN \S4.3). Its question: what fraction of
naive-grader successes survive your gate (\texttt{rexc trust-report}), by
cheat class?

\textbf{Fixed by the course.} The unit is a delivery: one run's final state.
The naive graders $g$ are five: the run's exit code; \texttt{cargo build} and
\texttt{cargo test}; a ``verified'' substring in the output; one verifier
seed; a run with \texttt{--allow-warnings}. For each, the inflation is $I_g =
1 - P(\text{gate PASS} \mid g \text{ PASS})$. The primary estimand is the
pooled $I$ over the deliveries that at least one naive grader passes,
estimated as a ratio of counts, with deliveries clustered by task. The primary
analysis is one estimate with a 95\% cluster-bootstrap interval (10,000
resamples of tasks) and no test; the inflation per cheat class and per naive
grader is descriptive, with Wilson intervals; the 4 impossible tasks are
reported separately. The data: every Lab $\alpha$ and Lab $\beta$ success,
re-graded with no new model call, plus 12 crate tasks (8 possible, 4
impossible), run by Claude (2 repeated runs) and by Qwen (3 repeats). Validity
rules: the crate tasks are built over \texttt{rexc}'s parser, character
classes, printer and CLI, never over the emitter, the simulation or jump
threading; \texttt{fvh} refuses to unpack \texttt{bench/local/r/crate/} until
your R16 statement files are committed (the registry's \texttt{statements}
pin, checked by \texttt{python -m fvh unseal R16}); CI never calls a model.

\textbf{Given numbers} (the course's budget for Lab $\gamma$, PLAN \S4.3).
The budget formula, per arm: calls $=$ tasks $\times$ attempts $\times$ rounds
$\times$ arms $\times$ repeats, plus the sentinel calls ($\Delta$148); runs
$=$ tasks $\times$ arms $\times$ repeats. Claude, harness-executed: 12 tasks
$\times$ 1 attempt $\times$ at most 25 rounds $\times$ 1 arm $\times$ 2
repeated runs, and at most 10 sentinel calls. Qwen in agent mode
(\texttt{opencode} in a per-run container): 12 tasks $\times$ 1 attempt
$\times$ at most 40 turns $\times$ 1 arm $\times$ 3 repeats. PLAN's working
figure is about 15 minutes of GPU time per Qwen run [A], to be replaced by a
measurement at your settings before the batch. The registry's R16 budget,
which \texttt{fvh} enforces: 610 Claude calls, 24 Claude runs, 1,440 Qwen
calls, 36 Qwen runs and 9 GPU-hours. PLAN's fallback if the measured cost puts
the plan over its ceiling: two Qwen repeats. Write the Qwen settings in the
terms of the apparatus's pinned OpenAI-compatible endpoint (the cap on output
tokens, the pinned context length), as for Lab $\beta$.

Write every field of the answer template: claim (and what would refute it);
metric and estimator (the unit; each naive grader's PASS rule; the gate's
PASS, FAIL and UNKNOWN and how UNKNOWN counts; $I_g$ and the pooled $I$ as a
ratio of counts); primary analysis (the interval's resampling written out,
with its seed; the per-class and per-grader Wilson intervals; the impossible
tasks); data (which Lab $\alpha$ and $\beta$ artifacts are re-graded, from
where and by what; the crate tasks; the arms, and the Claude runs as repeated
runs); budget arithmetic per arm, recomputed against the registry, with any
reconciliation stated; generation settings; cost per arm (tokens in and out,
GPU-seconds, wall clock); GPU-hours against the ceiling, including the
largest per-run cost your plan tolerates at a ceiling $C$; the truncation rate
as a sensitivity line; precision (there is no test: replace the MDE by the
expected half-width of the pooled $I$'s interval at $I \in \{0.1, 0.3, 0.5\}$
for about 40 to 80 naive successes, from a simulation you write, with its
assumptions and seed; pre-register the half-width, not a power; and say what
that precision cannot resolve); cheat surface and gate; missing results and
exclusions; \texttt{firsttry} intervals (80\% prediction intervals written
before any data); stopping rule; fallback; contamination; the statement-sha
unpack gate (what it checks, when, and what it protects); disclosure; and the
batch plan binding (\texttt{plan-sha256:}). Changing the block after grading
needs a new grade.
\end{problem}

\end{document}
